\chapter{Lorsque les routes se séparent}
\label{chap:routes-separent}

Les années avaient continué de passer au rythme des roues, des haltes et des
saisons.

Les deux enfants qui avaient autrefois approché la caravane sous les cris
d'alarme avaient grandi. À l'approche de leur vingtième année, ni Kael ni
Kamatsu n'avaient encore réellement quitté cette vie. La caravane restait leur
maison, celle qui les avait accueillis lorsqu'ils n'en avaient plus et dans
laquelle ils étaient devenus adultes.

Pourtant, depuis quelque temps, quelque chose changeait.

Kamatsu appréciait de plus en plus les villes où la caravane faisait étape. Il
pouvait y passer des heures à écouter les musiciens, observer les artisans,
découvrir les spectacles et discuter avec ceux qui avaient des histoires à
raconter. Ses carnets se remplissaient de dessins, de quelques paroles de
chansons, d'idées et de fragments de récits.

Il avait conservé tout ce que Rielle, Thom et Ellundril avaient éveillé en lui,
mais ces influences étaient désormais devenues les siennes.

Kael suivait une direction presque opposée.

Il aimait toujours les villes et leurs marchés, mais il était généralement
parmi les premiers à vouloir reprendre la route. Lors des haltes, il
s'éloignait de plus en plus souvent dans les bois, parfois avec son arc,
parfois simplement pour explorer.

Ses excursions de quelques heures étaient devenues des journées entières. Il
revenait avec du gibier lorsqu'il avait eu de la chance, quelques découvertes
lorsqu'il en avait eu moins, et presque toujours avec l'envie de repartir voir
ce qui se trouvait un peu plus loin.

Aucun des deux n'avait encore parlé de quitter la caravane.

Puis il y eut une fille.

Kael ne sut jamais exactement à quel moment elle avait pris suffisamment
d'importance pour modifier les projets de Kamatsu. Il savait seulement que,
lors de leurs dernières étapes dans une certaine ville, son frère trouvait
beaucoup plus souvent qu'autrefois d'excellentes raisons de s'y attarder.

Au début, Kael n'avait rien dit.

Puis il avait commencé à sourire chaque fois que Kamatsu annonçait qu'il devait
« aller voir quelque chose en ville ».

— Quoi ?

— Rien.

— Pourquoi tu souris ?

— Je souris pas.

— Si.

— Non.

Kamatsu avait fini par comprendre que nier ne servirait plus à grand-chose.

Lorsque la caravane commença à préparer son prochain départ, la plaisanterie
cessa pourtant d'en être une.

Cette fois, Kamatsu n'avait pas préparé ses affaires pour reprendre la route
avec les autres.

Il allait rester.

Pas uniquement pour elle, affirma-t-il.

Kael avait eu la délicatesse de ne pas trop insister sur ce point.

La ville lui plaisait. Il voulait découvrir autre chose, rencontrer des gens,
jouer, peindre, écrire peut-être. Il avait passé une grande partie de sa vie
sur les routes et, pour la première fois depuis longtemps, il avait envie de
voir ce qui se passerait s'il restait quelque part.

Kael comprenait.

Peut-être parce que, depuis quelque temps, il ressentait exactement l'inverse.

La route qu'il connaissait depuis des années ne lui suffisait plus tout à fait.
Robin lui avait appris à tirer, Artis avait éveillé sa curiosité pour les
soins, Cal lui avait appris à se déplacer et à observer, Ellundril lui avait
montré une autre manière de parcourir le monde et Durzo lui avait laissé
davantage de questions que de réponses.

À force de suivre les chemins tracés par les chariots, Kael avait commencé à
regarder ceux qui s'en éloignaient.

Il voulait savoir où ils menaient.

Il voulait partir dans les forêts, suivre des pistes dont il ne connaissait
pas la destination et découvrir ce qu'il était capable de faire lorsqu'aucun
chariot ne l'attendait le soir.

La décision de Kamatsu ne provoqua pas la sienne.

Mais, pour la première fois depuis leur rencontre, leurs routes allaient se
séparer.

\scenebreak

La veille du départ de la caravane, Kael parcourut une dernière fois le marché.

Il n'avait rien de particulier à acheter. Son équipement était prêt et ses
provisions suffisantes. Pourtant, devant l'étal d'un marchand, il s'arrêta.

Parmi divers objets se trouvaient plusieurs livres et carnets. Certains étaient
déjà remplis, d'autres presque inutilisables, mais l'un d'eux était vierge. Ce
n'était pas un ouvrage précieux : sa couverture était simple, le cuir portait
quelques imperfections et les pages n'étaient pas toutes parfaitement
régulières.

Kael le feuilleta.

Il pensa aux dessins de Kamatsu, aux chansons apprises auprès de Thom, aux
histoires racontées autour du feu et à toutes celles qu'ils avaient vécues sans
que personne les ait jamais écrites.

Il demanda le prix.

Comme autrefois pour son premier arc, Kael négocia assez peu et dépensa
probablement plus qu'il ne l'aurait souhaité.

Le lendemain matin, alors que les derniers préparatifs occupaient la caravane,
il retrouva Kamatsu et lui tendit le livre.

— Tiens.

Kamatsu le prit, surpris.

— C'est quoi ?

— Un livre.

— J'avais remarqué.

— Alors pourquoi tu demandes ?

Kamatsu leva les yeux au ciel avant de l'ouvrir. Il tourna quelques pages et
comprit qu'elles étaient toutes vierges.

— Il n'y a rien dedans.

— Justement.

Kamatsu releva la tête.

Kael haussa les épaules, soudain beaucoup moins à l'aise qu'il ne l'avait été
lorsqu'il l'avait acheté.

— Thom disait qu'il fallait pas passer sa vie à raconter les histoires des
autres.

Kamatsu regarda de nouveau les pages.

— Alors je me suis dit que tu pourrais écrire les tiennes.

Cette fois, Kamatsu ne trouva immédiatement aucune plaisanterie à faire.

Kael ajouta :

— Et puis, comme ça, la prochaine fois qu'on se verra, tu auras des histoires à
raconter.

Un sourire apparut finalement sur le visage de son frère.

— Les histoires de Kamatsu ?

— Si tu veux les appeler comme ça.

Kamatsu passa lentement une main sur la couverture.

— J'aime bien.

Il referma soigneusement le livre.

— Attends.

Il partit fouiller dans ses affaires et revint quelques instants plus tard avec
quelque chose enveloppé dans un morceau de tissu.

— J'avais prévu quelque chose aussi.

Kael prit le paquet.

À l'intérieur se trouvait une dague.

Il la sortit de son fourreau et l'examina. Ce n'était pas une arme d'apparat.
Elle avait été faite pour servir : une lame solide, une poignée simple,
suffisamment compacte pour rester facilement accessible sur la route.

Kael leva les yeux vers Kamatsu.

— Pourquoi une dague ?

— Parce qu'un livre, dans la forêt, ça ne va probablement pas beaucoup
t'aider.

Kael sourit.

— Tu pourrais frapper quelqu'un avec.

— Je préfère te donner ça.

Kael observa encore l'arme avant de la remettre dans son fourreau.

La première arme qu'il avait réellement appris à manier dans son village avait
été une dague. Des années plus tard, son frère lui en offrait une au moment
précis où leurs routes allaient se séparer.

Il ne savait pas si Kamatsu avait pensé à ce symbole.

Il ne lui demanda pas.

\scenebreak

Autour d'eux, la caravane terminait ses préparatifs. Les chevaux étaient
attelés, les dernières caisses chargées et les voix familières se répondaient
d'un chariot à l'autre.

Cette agitation avait accompagné une grande partie de leur vie.

Cette fois pourtant, pour la première fois, ils n'allaient pas repartir
ensemble.

Kamatsu resterait dans la ville.

Kael, lui, reprendrait la route avec la caravane.

Pour combien de temps, il n'en savait rien.

Depuis plusieurs mois déjà, il sentait grandir cette envie de s'éloigner des
chemins tracés par les chariots, de suivre une piste simplement pour découvrir
où elle menait et de savoir ce dont il serait capable lorsqu'aucun campement
familier ne l'attendrait le soir.

Le départ de Kamatsu ne provoquait pas encore le sien, mais il rendait une
chose évidente : tôt ou tard, Kael partirait lui aussi.

Ils avaient passé tant d'années ensemble que le moment où ils durent réellement
se séparer leur parut étrangement difficile à remplir de mots.

Kamatsu regarda l'équipement de son frère.

— Tu sais où va la caravane après ?

Kael acquiesça et lui indiqua leur prochaine étape.

Kamatsu eut un petit sourire.

— Et après ?

Kael haussa les épaules.

— Après, on verra.

— Ça te ressemble déjà davantage.

Kael lui donna un léger coup d'épaule.

Pendant un instant, ils retrouvèrent exactement la complicité que cet inconnu
avait remarquée des années auparavant.

Puis Kamatsu désigna la dague.

— Essaie de ne pas la perdre.

— Essaie de remplir le livre.

— Ça risque d'être long.

— T'as le temps.

Ils se regardèrent.

Il n'y eut pas de grande promesse. Ils ne jurèrent pas de se retrouver à une
date précise ni de revenir au même endroit chaque année. Leur vie s'était
construite sur les routes ; tous deux savaient mieux que beaucoup que les
promesses de ce genre résistaient mal à la distance et au hasard.

Mais ils savaient aussi autre chose.

Une route pouvait se séparer en deux sans effacer celle qui avait précédé.

Kael serra son frère contre lui. Kamatsu lui rendit son étreinte et, pendant
quelques secondes, ils redevinrent presque les deux garçons qui avaient fui
ensemble des années auparavant.

Lorsqu'ils se séparèrent, Kamatsu avait toujours le livre à la main.

Kael portait désormais la dague à sa ceinture.

Ils avaient autrefois tout perdu presque au même moment, puis trouvé ensemble
une nouvelle famille sur les routes. Ils avaient grandi dans les mêmes
chariots, autour des mêmes feux, appris auprès des mêmes voyageurs et découvert
peu à peu des chemins qui leur appartenaient.

Kamatsu avait choisi de commencer à suivre le sien.

Kael savait que son tour viendrait bientôt.

Lorsque les premiers chariots commencèrent à s'ébranler, Kael rejoignit la
caravane tandis que Kamatsu resta en arrière.

Avant que la distance ne devienne trop grande, Kael se retourna.

Kamatsu avait fait exactement la même chose.

Ils échangèrent un dernier signe.

Puis la caravane poursuivit sa route.

Ce jour-là, leurs chemins se séparèrent pour la première fois depuis leur
rencontre.

Mais ils ne cessèrent pas d'être frères.